\section*{Chapitre \ref{ch:mx:tick-size-queues-and-priority} --- Pas de cotation, files d'attente et priorité}

\begin{solution}{exo:mx:tick-size-queues-and-priority:1}
$0.01/4=25$ points de base et $0.01/400=0.25$ points de base. L'action à quatre dollars a plus de chances d'être un \omterm{def:mx:tick-size-queues-and-priority:large}{actif à gros tick}, mais la réponse dépend de la fourchette
qu'elle aurait sans la contrainte, qui dépend de sa volatilité et de son activité~: une action très active et peu volatile peut être un \omterm{def:mx:tick-size-queues-and-priority:large}{actif à gros tick} à un prix élevé.
\end{solution}

\begin{solution}{exo:mx:tick-size-queues-and-priority:2}
$\hat\eta=40/(2\times160)=0.125$~; fourchette implicite $2\times0.125\times1=0.25$ cents. Oui~: $\eta$ est nettement inférieur à un demi.
\end{solution}

\begin{solution}{exo:mx:tick-size-queues-and-priority:3}
$P_2=\tfrac{0.2}{0.25}\cdot\tfrac{0.25}{0.30}\cdot\tfrac{0.30}{0.35}=0.571$. Temps moyen~: $1/0.2+1/0.25+1/0.3=12.33$ secondes.
\end{solution}

\begin{solution}{exo:mx:tick-size-queues-and-priority:4}
Tête de file~: $0.7\times(2.5-0.4)=1.47$ cents par action~; fin de file~: $0.06\times(2.5-0.9)=0.096$. La tête de file vaut 1.37 cents par action de plus~: 6.87 dollars pour 500
actions.
\end{solution}

\begin{solution}{exo:mx:tick-size-queues-and-priority:5}
Avant--après~: $9\,960-2\,400=7\,560$. \omterm{def:mx:tick-size-queues-and-priority:did}{Doubles différences}~: $7\,560-(2\,350-2\,750)=7\,960$. C'est causal si, sans la nouvelle grille, la
profondeur des actions traitées aurait évolué comme celle des témoins (tendances parallèles)~; c'est le cas ici par construction, car le simulateur tire les actions de façon identique.
\end{solution}

\begin{solution}{exo:mx:tick-size-queues-and-priority:6}
Les trois fenêtres d'une même action partagent sa liquidité et ses tirages aléatoires~; elles sont donc positivement corrélées et portent moins d'information que trois
observations indépendantes~; la formule qui suppose l'indépendance l'ignore et sous-estime l'erreur type d'un facteur 1.7. Il faut retenir les erreurs types groupées~; avec seulement 16
groupes, une correction pour petit échantillon ou un bootstrap wild est encore plus prudent.
\end{solution}

\begin{solution}{exo:mx:tick-size-queues-and-priority:7}
Deux heures par règle, avec les graines 7 et 8~: sous priorité temporelle, 0.70 des ordres en tête de file et 0.05 de ceux placés derrière 20 ordres ou plus reçoivent une exécution, valant 1.28 et 0.081 cents
par action affichée~; sous pro rata, 0.78 et 0.58, valant 0.72 et 0.093. Le pro rata attribue à chaque ordre au repos une part de chaque exécution proportionnelle à sa taille,
si bien qu'un ordre situé loin dans une longue file ne reçoit encore qu'une petite part de sa taille, et la valeur par action affichée baisse avec la longueur de la file qu'elle
partage.
\end{solution}

\begin{solution}{exo:mx:tick-size-queues-and-priority:8}
La profondeur au meilleur cours augmente mécaniquement lorsque plusieurs niveaux à un cent sont regroupés à un prix de cinq cents~: les mêmes actions sont comptées à moins de prix différents. Le
services de la SEC n'ont constaté aucune variation significative de la profondeur à moins de cinq cents du meilleur cours pour les deux premiers groupes, et les \omterm{def:mx:decomposing-the-spread:quoted}{fourchettes affichées} ont augmenté de cinq à sept points de base
contre le groupe témoin. La simulation donne le même résultat~: la profondeur au meilleur cours est multipliée par 3.8 sur une grille à cinq cents sans aucun nouvel ordre
du tout, tandis que tout \omterm{def:mx:the-limit-order-book:orders}{ordre au marché} paie 4.5 fois plus.
\end{solution}

\begin{solution}{pb:mx:tick-size-queues-and-priority:1}
\textbf{1.} Une action dont la fourchette vaut presque toujours un tick~; il faut comparer le tick à la fourchette que l'actif aurait sans contrainte, ce qui dépend de
sa volatilité et de son activité.
\textbf{2.} Il maintient fixe le flux d'ordres de \texttt{firm.tape} et ne change que la grille de prix (prix limites arrondis en s'éloignant du marché).
\textbf{3.} 1.07 cents et 94.2~\%~; 5.03 cents et 99.4~\%.
\textbf{4.} 1\,740 à 6\,560 actions, 8.7 à 33.6 ordres.
\textbf{5.} 0.56 et 2.53 cents de demi-fourchette effective~: les mêmes \omterm{def:mx:the-limit-order-book:orders}{ordres au marché} traversent une fourchette qui vaut désormais un tick de cinq cents.
\textbf{6.} Profit espéré par action en rejoignant la file~; $V(n)=P_n(c-a(n))$.
\textbf{7.} 11.3, 24.1, 55.7 et 81.4 secondes.
\textbf{8.} 0.76, 0.44 et 0.13 prédits pour 0, 3 et 20 ordres devant, contre 0.54, 0.22 et 0.06~: le modèle ignore les nouveaux niveaux placés devant et des annulations plus rapides
lorsque le prix se déplace.
\textbf{9.} Tête de file 0.57, 0.66, 1.26 et 2.05 cents par action affichée~; fin de file 0.001, 0.005, 0.094 et 0.188.
\textbf{10.} La fin de file ne s'exécute que lorsque le flux a consommé toute la file, ce qui arrive plutôt avec un flux qui continue de déplacer le prix~: à un cent, $-0.71$ cents
en tête de file contre 0.48 derrière 20 ordres.
\textbf{11.} $\hat\eta=N_c/(2N_a)$ sur les variations d'un tick du prix négocié.
\textbf{12.} 0.108, 0.058, 0.019 et 0.007~; 0.22, 0.23, 0.19 et 0.14 cents.
\textbf{13.} La fourchette implicite est proche d'une propriété de l'actif~: la \omterm{def:mx:decomposing-the-spread:quoted}{fourchette affichée} suit la grille, alors que le coût propre à l'actif d'une variation de prix ne la suit pas.
\textbf{14.} Les actions NMS cotées à 1.00 dollar ou plus dont la \omterm{def:mx:decomposing-the-spread:quoted}{fourchette affichée} moyenne pondérée dans le temps vaut 0.015 dollar ou moins, pour les cotations et les ordres.
\textbf{15.} Environ 400 actions chacune~: cotation au pas de cinq cents~; cotation et négociation au pas de cinq cents~; la même chose, plus trade-at.
\textbf{16.} Les fourchettes ont augmenté dans chaque groupe entre les deux périodes, groupe témoin compris.
\textbf{17.} \omterm{def:mx:decomposing-the-spread:quoted}{Fourchettes affichées} en hausse de cinq à sept points de base par rapport au groupe témoin~; profondeur au meilleur cours en hausse, surtout mécaniquement~; profondeur à moins de cinq cents sans
changement significatif dans les deux premiers groupes.
\textbf{18.} $+7\,960$ actions~; 2\,270 avec erreurs types groupées par action, 1\,340 sans groupement.
\textbf{19.} \emph{Résultat nommé}~: multiplier le tick par cinq sur le même flux multiplie la fourchette par 4.7 (de 1.07 à 5.03 cents), la taille affichée au
meilleur cours par 3.8, le temps médian avant exécution par 4.9 et la valeur de la tête de file par 2.2 (de 0.57 à 1.26 cents par action affichée)~; le pilote a mesuré une
hausse de cinq à sept points de base des \omterm{def:mx:decomposing-the-spread:quoted}{fourchettes affichées} par rapport aux témoins et une hausse surtout mécanique de la profondeur au meilleur cours.
\textbf{20.} Les fournisseurs de liquidité en tête de file, rémunérés par chaque ordre qui traverse la fourchette la plus large.
\end{solution}

\begin{solution}{iq:mx:tick-size-queues-and-priority:1}
Sa fourchette vaut presque toujours un tick et ses files sont longues~; la négocier relève de la position dans la file et du timing (rejoindre tôt, sans payer la
fourchette), alors qu'une action à petit tick a un carnet clairsemé, des variations de fourchette et une concurrence par les prix.
\iqlookfor{Fourchette bloquée à un tick~; files d'attente~; concurrence prix contre temps.}
\end{solution}

\begin{solution}{iq:mx:tick-size-queues-and-priority:2}
La tête de file s'exécute plus souvent et face à un flux moins informé~; la fin de file s'exécute rarement, surtout lorsque le prix la traverse. Elle vaut davantage lorsque le
tick est grand par rapport à la fourchette naturelle de l'actif, que les files sont longues et que la priorité est temporelle.
\iqlookfor{Les deux facteurs (probabilité et sélection adverse conditionnelle)~; le tick~; la règle de priorité.}
\end{solution}

\begin{solution}{iq:mx:tick-size-queues-and-priority:3}
Une file de naissance et de mort~: les \omterm{def:mx:the-limit-order-book:orders}{ordres au marché} au taux $\mu$ prennent la tête, les ordres devant annulent au taux $\theta$, les nôtres au taux $\nu$~: $P_n=\tfrac{\mu}{\mu+\nu}\prod_k\tfrac{\mu+k
\theta}{\mu+k\theta+\nu}$. Il ignore les variations de prix (apparition d'un meilleur niveau, vidage du niveau), les annulations dépendant de l'état, les tailles des \omterm{def:mx:the-limit-order-book:orders}{ordres au marché} et les
\omterm{def:mx:order-types-and-their-uses:hidden}{ordres cachés}.
\iqlookfor{La formule~; la dynamique des prix qui lui manque.}
\end{solution}

\begin{solution}{iq:mx:tick-size-queues-and-priority:4}
\omterm{def:mx:tick-size-queues-and-priority:did}{Doubles différences}~: régresser les fourchettes sur le groupe traité, sur la période postérieure et sur leur interaction, avec contrôles par action et par date, erreurs types groupées par action~;
vérifier les tendances préalables et l'équilibre des groupes, et présenter les résultats par fourchette de la période préalable (la contrainte ne joue que là où l'ancienne fourchette était inférieure au nouveau tick).
\iqlookfor{Groupe témoin~; interaction~; groupement~; hétérogénéité.}
\end{solution}

\begin{solution}{iq:mx:tick-size-queues-and-priority:5}
La position compte moins et la taille davantage~: les exécutions sont réparties selon la taille, si bien que j'affiche plus que ce que je veux négocier, et je gère le risque de surexécution. Je surveille les règles du top order et
de l'allocation minimale, et j'annule plus vite lorsque la taille totale de la file pourrait trop me remplir d'un coup.
\iqlookfor{Surdimensionnement~; règles d'allocation~; risque d'exécution.}
\end{solution}

\begin{solution}{iq:mx:tick-size-queues-and-priority:6}
Il faut vérifier la part du temps où la fourchette vaut un tick et la longueur des files, puis estimer $\eta$ à partir des continuations et des alternances à un tick~: $\eta$ nettement inférieur
à un demi, avec une fourchette bloquée à un tick, signifie que le tick est grand par rapport à la fourchette implicite de l'actif $2\eta\delta_{\mathrm{tick}}$.
\iqlookfor{Part à un tick~; longueurs de file~; estimateur de la zone d'incertitude.}
\end{solution}
